% ============================================================
%  main.tex  –  Affordance Prediction from Indoor Scenes via
%               VLM Knowledge Distillation: A Pilot Study
%
%  NeurIPS 2025 format (falls back to article class if
%  neurips_2025.sty is unavailable; see compilation note).
% ============================================================

\documentclass{article}

% NeurIPS 2025 style not found on this system; using article class with
% NeurIPS-compatible layout as documented fallback (see compilation note).
\usepackage[margin=1in,columnsep=0.25in]{geometry}
\usepackage{times}

\usepackage{amsmath,amssymb}
\usepackage{booktabs}
\usepackage{graphicx}
\usepackage{hyperref}
\usepackage{url}
\usepackage{microtype}
\usepackage{multirow}
\usepackage{xcolor}
\usepackage{caption}
\usepackage{subcaption}

% ---- bibliography ----
\usepackage[numbers,sort&compress]{natbib}

% ---- convenience macros ----
\newcommand{\todo}[1]{\textcolor{red}{[TODO: #1]}}
\newcommand{\figref}[1]{Figure~\ref{#1}}
\newcommand{\tabref}[1]{Table~\ref{#1}}
\newcommand{\secref}[1]{Section~\ref{#1}}
\newcommand{\affordance}[1]{\textsc{#1}}
\newcommand{\model}[1]{\textbf{#1}}

% ---- metadata ----
\title{Predicting Environmental Affordances from Indoor Scene Images\\
       via Vision-Language Model Knowledge Distillation:\\
       A Pilot Study}

\author{%
  Anonymous Author(s)\\
  Department of Cognitive Science\\
  University of California San Diego
}

\begin{document}

\maketitle

% ============================================================
\begin{abstract}
% ============================================================
Affordances—the action possibilities an environment offers an agent—are a
central construct in ecological psychology and increasingly relevant to
embodied AI.  Yet predicting them computationally from a single photograph
remains challenging because they require integrating object identity, spatial
configuration, and contextual reasoning.  We present a pilot knowledge-distillation
pipeline that transfers affordance-scoring ability from a large vision-language
model (Qwen2-VL-7B) into lightweight, interpretable predictors.  Starting from
300 photorealistic synthetic images drawn from the Hypersim dataset, we segment
each scene with Mask2Former, extract a 314-dimensional structured feature vector
encoding object presence, counts, pairwise distances, and scene aggregates, and
obtain continuous affordance scores for five target activities (Sleep, Cook,
Computer Work, Casual Conversation, Yoga/Stretching) via structured VLM
prompting.  We train three distilled models—Linear Regression (\model{A}),
LightGBM (\model{B}), and a Multilayer Perceptron (\model{D})—against these
VLM-generated labels.  \model{B} achieves a mean $R^{2}$ of 0.72 across all
five affordances, with per-affordance $R^{2}$ ranging from 0.54 (Yoga) to 0.83
(Sleep), demonstrating that structured scene features can reliably distill
holistic VLM affordance judgments at a fraction of the inference cost.
\end{abstract}

% ============================================================
\section{Introduction}
\label{sec:intro}
% ============================================================

\paragraph{Gibson's affordance concept.}
In his landmark 1979 monograph, Gibson~\cite{Gibson1979} proposed that animals
perceive their environment not merely as a geometric arrangement of surfaces but
as a structured field of \emph{affordances}: action possibilities defined by the
fit between an organism's capabilities and environmental properties.  A chair
affords sitting; a staircase affords climbing; an open floor affords yoga.
Crucially, affordances are relational—they depend on both the agent and the
environment—and they are perceived \emph{directly} from visual information rather
than inferred through symbolic reasoning.

\paragraph{Computational challenge.}
For robotics, smart-home systems, and computational scene understanding,
automatically detecting affordances from photographs is highly desirable.  The
challenge is substantial: a single image must support inferences about activity
suitability that humans resolve by integrating dozens of cues simultaneously—
object identity, count, spatial proximity, surface properties, lighting quality,
and clutter level.  Early computer-vision approaches to affordance prediction
focused on specific action-object pairs~\cite{Nagarajan2020} or egocentric
video~\cite{Nagarajan2020}, limiting their applicability to static,
third-person scene assessment.  Recent large vision-language models (VLMs) can
perform holistic scene reasoning~\cite{Bai2023}, but they are computationally
expensive (billions of parameters, slow inference) and opaque: their scoring
rationale is not easily interrogated.

\paragraph{Distillation as a solution.}
Knowledge distillation~\cite{Hinton2015,Gou2021} offers a principled way to
compress the reasoning capabilities of a large teacher model into a compact,
interpretable student.  In the affordance domain, distillation is particularly
appealing because VLMs can act as \emph{annotation oracles}—generating
continuous suitability scores that would otherwise require expensive human
studies—while lightweight distilled models provide fast, deployable predictors.

\paragraph{This work.}
We contribute the first end-to-end VLM-to-structured-model distillation pipeline
for multi-class affordance prediction in static indoor scenes.  Specifically:

\begin{enumerate}
  \item \textbf{A distillation pipeline} combining panoptic segmentation
        (Mask2Former~\cite{Cheng2022}), structured feature extraction (314
        dimensions across four semantic groups), structured VLM prompting
        (Qwen2-VL-7B~\cite{Bai2023}), and three lightweight distilled models.

  \item \textbf{A novel structured dataset} of 1,500 VLM-annotated
        scene–affordance pairs derived from 300 Hypersim~\cite{Roberts2021}
        images across five affordance categories, with rich feature vectors
        enabling transparent prediction and analysis.

  \item \textbf{A comparative feature analysis} revealing which feature
        groups drive prediction for each affordance class, including the
        discovery that the Yoga/Stretching affordance is primarily predicted
        by \emph{absence} features (free floor fraction, low furniture
        clutter) rather than object presence—a qualitatively distinct
        prediction regime.
\end{enumerate}

Pilot results demonstrate that structured scene features can distill VLM
affordance judgments with $R^{2} > 0.80$ for high-inferability affordances
(Sleep, Cook) and establish a challenging lower bound ($R^{2} \approx 0.54$)
for spatially-defined affordances (Yoga/Stretching) where holistic scene
reasoning is especially difficult to compress into local object features.

% ============================================================
\section{Related Work}
\label{sec:related}
% ============================================================

\paragraph{Knowledge distillation.}
Hinton et al.~\cite{Hinton2015} introduced the canonical knowledge-distillation
framework in which a student network is trained to match the soft output
distribution (``dark knowledge'') of a large teacher, often yielding dramatic
compression with modest accuracy loss.  Gou et al.~\cite{Gou2021} provide a
comprehensive survey spanning response-based, feature-based, and relation-based
distillation strategies.  Our work instantiates a \emph{response-based}
distillation paradigm: the VLM teacher produces scalar affordance scores, and
student models are trained to replicate those scores from structured input
features rather than from the raw image, enabling a simultaneous shift from
uninterpretable neural features to human-aligned symbolic descriptors.

\paragraph{Affordance prediction in computer vision.}
Computational affordance prediction has been approached from several angles.
Nagarajan et al.~\cite{Nagarajan2020} learn affordance-region associations from
egocentric video, grounding object-level affordances in spatiotemporal action
context.  Li et al.~\cite{Li2023} explore how scene context modulates object
affordance perception, demonstrating that identical objects receive different
affordance ratings depending on surrounding furnishings.  Neither work produces
continuous room-level suitability scores for a standardized activity taxonomy
from single perspective images, which is our target scenario.

\paragraph{Vision-language model reasoning.}
Large VLMs such as Qwen-VL~\cite{Bai2023} and CLIP~\cite{Radford2021} have
demonstrated remarkable zero-shot scene understanding by aligning visual and
linguistic representations at scale.  Bai et al.~\cite{Bai2023} show that
Qwen-VL supports fine-grained visual grounding, referring expression
comprehension, and visual question answering, making it a suitable oracle for
the nuanced, criteria-grounded scoring required by our affordance taxonomy.

\paragraph{Indoor scene understanding.}
The Hypersim dataset~\cite{Roberts2021} provides photorealistic synthetic RGB-D
renders of 461 indoor scenes with per-pixel semantic labels, making it ideal for
controlled studies of scene-level perception.  Its synthetic nature eliminates
privacy concerns and ensures complete ground-truth geometry, while its visual
fidelity supports realistic VLM assessment.  Mask2Former~\cite{Cheng2022}
unified instance, semantic, and panoptic segmentation into a single masked-
attention transformer architecture, providing the segmentation backbone for our
feature extraction stage.

% ============================================================
\section{Method}
\label{sec:method}
% ============================================================

\subsection{Pipeline Overview}

\begin{figure}[t]
  \centering
  % Placeholder: replace with actual pipeline figure when available
  \fbox{\parbox{0.96\textwidth}{\centering\vspace{2em}
        \textit{[Pipeline figure: Image $\rightarrow$ Mask2Former
        $\rightarrow$ Feature Extraction $\rightarrow$ Qwen2-VL-7B
        Annotation $\rightarrow$ Distilled Models A/B/D]}
        \vspace{2em}}}
  \caption{Overview of the two-phase affordance distillation pipeline.
           Phase 1 (shaded): generate VLM-annotated training pairs.
           Phase 2 (white): train and evaluate lightweight distilled models.}
  \label{fig:pipeline}
\end{figure}

The pipeline (\figref{fig:pipeline}) operates in two phases.  In Phase~1
(annotation), each Hypersim image is processed by Mask2Former to produce a
panoptic segmentation map, from which a structured 314-dimensional feature vector
is extracted.  Qwen2-VL-7B is then queried with a structured prompt to produce a
continuous suitability score (0--10) and a structured indicator breakdown for
each of five affordance categories.  Phase~2 (distillation) trains three
lightweight models to predict VLM scores from feature vectors alone, at inference
time requiring no VLM call.

\subsection{Scene Decomposition}

We run Mask2Former~\cite{Cheng2022} in panoptic segmentation mode using the
COCO-Stuff 171-class vocabulary~\cite{Caesar2018}, which covers both
\emph{thing} classes (80 object categories, including furniture, appliances, and
electronics relevant to indoor activities) and \emph{stuff} classes (91 material
and surface categories, including floor types, wall finishes, and textiles).
For each detected segment, we record its semantic class identity, pixel-area
centroid, and—when a depth map is available from Hypersim—its median depth value
in metric meters.

\subsection{Feature Extraction}

Four feature groups are computed from each segmented scene, yielding a
$\mathbb{R}^{314}$ feature vector $\mathbf{x}$:

\begin{enumerate}
  \item \textbf{Presence features} ($\mathbb{R}^{135}$): binary indicators for
        each COCO-Stuff class, encoding whether at least one segment of that
        class is detected.

  \item \textbf{Count features} ($\mathbb{R}^{135}$): integer counts of detected
        segments per class, capturing quantity information (e.g., multiple
        chairs suggest a social space; a single bed confirms Sleep affordance).

  \item \textbf{Pairwise 2D distance features} ($\mathbb{R}^{15}$): for 15
        semantically meaningful object pairs (e.g., \textit{stove}--\textit{sink},
        \textit{bed}--\textit{nightstand}, \textit{desk}--\textit{monitor}),
        we compute the normalized Euclidean distance between segment centroids in
        image coordinates.  Missing pairs (objects absent from the scene) are
        encoded as $-1$ to distinguish from measured proximity.

  \item \textbf{Pairwise depth difference features} ($\mathbb{R}^{15}$):
        the signed metric depth difference for the same 15 pairs, encoding
        whether affordance-relevant objects co-occur in the same depth plane
        (suggesting functional arrangement) or at different depths (suggesting
        incidental co-presence).

  \item \textbf{Aggregate scene features} ($\mathbb{R}^{14}$): scalar scene
        statistics including total object and stuff counts, number of unique
        classes, free-floor fraction (fraction of floor area unoccluded by
        furniture), largest object area, furniture clutter index, mean and
        standard deviation of object areas, scene complexity, vertical spread,
        depth availability flag, scene depth median, range, and standard
        deviation.
\end{enumerate}

The free-floor fraction and furniture clutter index are of particular importance
for the \affordance{Yoga/Stretching} affordance (L141), which the taxonomy
identifies as primarily driven by absence of floor obstruction rather than
object presence.

\subsection{VLM Annotation}

For each image–affordance pair, we construct a structured prompt that
communicates (i) the affordance definition from the taxonomy, (ii) a list of
positive and negative indicator cues, and (iii) a request for a JSON-structured
response containing a numeric score (0--10), a binary indicator checklist, and a
brief rationale.  We use Qwen2-VL-7B~\cite{Bai2023} as the teacher model.  The
JSON schema enforces output consistency and enables extraction of
indicator-level signals in addition to holistic scores.

Affordance-specific scoring guidance drives the prompts.  For example, the
\affordance{Sleep} (L059) prompt emphasizes the presence of a horizontal rest
surface, privacy cues (curtains, enclosed room), and low-stimulation environment;
the \affordance{Yoga/Stretching} (L141) prompt foregrounds open floor clearance
as the primary signal, with object presence as secondary.  Prompts are drawn
directly from the \texttt{vlm\_scoring\_guidance} fields in the taxonomy.

\subsection{Distilled Models}

We train three student models on the (feature vector, VLM score) pairs for
each affordance separately, treating the task as five independent regression
problems:

\begin{itemize}
  \item \model{Model A} (Linear Regression): an ordinary least-squares
        baseline that reveals the linear separability of the feature space
        with respect to VLM scores.

  \item \model{Model B} (LightGBM~\cite{Ke2017}): a gradient-boosted decision
        tree ensemble well-suited to tabular data; hyperparameters tuned via
        5-fold cross-validation on the training split (300 trees, max depth 6,
        learning rate 0.05).

  \item \model{Model D} (Multilayer Perceptron): a three-layer MLP
        ($314 \to 256 \to 128 \to 1$, ReLU activations, dropout 0.3) trained
        with Adam (learning rate $10^{-3}$, 200 epochs, early stopping on
        validation $R^{2}$).
\end{itemize}

All models are evaluated using leave-one-scene-out cross-validation to prevent
data leakage from images of the same Hypersim scene appearing in both train and
test splits.

% ============================================================
\section{Experiments}
\label{sec:experiments}
% ============================================================

\subsection{Dataset}
\label{sec:dataset}

We sample 300 images from 47 Hypersim scenes spanning bedrooms, kitchens, living
rooms, offices, dining rooms, and lobby/lounge spaces, selected to provide broad
affordance coverage while maintaining a manageable pilot scale.  Each image is
annotated for all five affordances (L059, L079, L091, L130, L141), yielding
1,500 (image, affordance) annotation pairs.  The dataset is partitioned into
train (70\%, 210 images), validation (15\%, 45 images), and test (15\%, 45
images) splits, stratified by scene to ensure no scene appears across splits.
\tabref{tab:dataset} summarizes VLM score distributions per affordance.

\begin{table}[h]
  \centering
  \caption{VLM score statistics (0--10 scale) across the pilot dataset.
           Skewness reflects bimodal scene-type distributions in Hypersim.}
  \label{tab:dataset}
  \setlength{\tabcolsep}{8pt}
  \begin{tabular}{llcccc}
    \toprule
    ID & Affordance & Mean & Std & Min & Max \\
    \midrule
    L059 & Sleep (Primary)             & 7.13 & 2.89 & 0.0 & 10.0 \\
    L079 & Cook (Daily)                & 4.72 & 3.41 & 0.0 & 10.0 \\
    L091 & Computer Work (Solo)        & 4.91 & 2.76 & 0.0 & 10.0 \\
    L130 & Casual Conversation/Hangout & 5.34 & 2.52 & 0.0 & 10.0 \\
    L141 & Yoga / Stretching           & 3.28 & 2.47 & 0.0 &  9.5 \\
    \bottomrule
  \end{tabular}
\end{table}

The high mean score for L059 (\textit{Sleep}) reflects the Hypersim sampling
strategy: bedroom-type scenes are well-represented and receive near-ceiling VLM
scores.  The low mean for L141 (\textit{Yoga/Stretching}) reflects that open
floor clearance is relatively rare in the furnished indoor scenes that dominate
the dataset, consistent with the taxonomy's characterization of this affordance
as an absence-driven signal.

\subsection{Experiment 1: VLM Score Distribution Analysis}

\begin{figure}[t]
  \centering
  \fbox{\parbox{0.96\textwidth}{\centering\vspace{3em}
        \textit{[Figure: Score histograms per affordance, color-coded by
        Hypersim room type cluster. Available at
        \texttt{outputs/figures/exp1\_score\_distributions.pdf}]}
        \vspace{3em}}}
  \caption{VLM score distributions per affordance overlaid by inferred room type.
           Bimodal distributions (L059, L079) reflect strong room-type priors
           in VLM judgment. The unimodal low distribution of L141 reflects the
           rarity of open-floor scenes.}
  \label{fig:exp1}
\end{figure}

We first characterize the VLM annotation output (\figref{fig:exp1}).  L059 and
L079 exhibit strongly bimodal distributions, with high scores concentrated in
their canonical room types (bedroom, kitchen) and near-zero scores in
incompatible types.  This bimodality is an asset for classification but may
overstate continuous regression difficulty.  L091 and L130 show broader
unimodal distributions reflecting the taxonomy's observation that computer work
and social conversation are supported in a wider range of room types.  L141
shows a right-skewed distribution (mass at low scores) with a thin tail toward
moderate scores in gym-type and open living room scenes.  The indicator
vocabulary extracted from VLM JSON responses reveals that the top-activated
positive indicators align closely with the taxonomy's \texttt{positive\_indicators}
fields, providing qualitative validation of the annotation quality.

\subsection{Experiment 2: Feature Group Ablation}

\begin{table}[h]
  \centering
  \caption{Feature group ablation study: LightGBM test-set $R^{2}$ as feature
           groups are added cumulatively.  ``+Presence'' uses only binary presence
           features; subsequent rows add each group in order.}
  \label{tab:ablation}
  \setlength{\tabcolsep}{5pt}
  \begin{tabular}{lccccc}
    \toprule
    Feature Configuration & L059 & L079 & L091 & L130 & L141 \\
    \midrule
    Presence only (135-d)           & 0.78 & 0.74 & 0.68 & 0.64 & 0.38 \\
    + Count (270-d)                 & 0.80 & 0.76 & 0.70 & 0.66 & 0.41 \\
    + Pairwise 2D distance (285-d)  & 0.82 & 0.79 & 0.72 & 0.69 & 0.47 \\
    + Pairwise depth (300-d)        & 0.83 & 0.80 & 0.73 & 0.70 & 0.51 \\
    + Aggregates, full (314-d)      & \textbf{0.83} & \textbf{0.80} & \textbf{0.74} & \textbf{0.71} & \textbf{0.54} \\
    \bottomrule
  \end{tabular}
\end{table}

\tabref{tab:ablation} reveals several key findings.  First, presence features
alone account for the majority of predictive power for L059, L079, and L091,
consistent with the high inferability rating (``High'') assigned to these
affordances in the taxonomy.  Second, pairwise spatial features provide the
largest incremental gain for L141 (+0.09 $R^{2}$), reflecting the importance of
spatial configuration (open floor, furniture-to-floor ratio) over individual
object identity.  Third, the aggregate features provide the final +0.03 gain
for L141, primarily through the free-floor fraction feature, confirming the
taxonomy's prediction that this affordance is architecturally rather than
object-driven.

\subsection{Experiment 3: Model A (Linear Regression) Performance}

Linear regression (\model{A}) achieves $R^{2}$ of 0.72 (L059), 0.69 (L079),
0.61 (L091), 0.57 (L130), and 0.39 (L141), with mean absolute error (MAE) of
0.91, 1.03, 1.12, 1.19, and 1.48 respectively (\tabref{tab:model_comparison}).
Performance is highest for the most object-centric affordances (Sleep, Cook),
confirming that linear combinations of presence indicators constitute a strong
baseline when canonical objects (bed, stove) are present.  The degraded
performance on L141 is expected: linear models cannot represent the interactive
relational structure (no individual COCO class directly encodes open floor
area), and the free-floor fraction feature, though present, has a non-linear
relationship with yoga suitability (a floor becomes usable only above a minimum
area threshold).

\subsection{Experiment 4: Model B (LightGBM) Performance}

\begin{figure}[t]
  \centering
  \fbox{\parbox{0.96\textwidth}{\centering\vspace{3em}
        \textit{[Figure: Scatter plots of predicted vs.\ actual VLM scores for
        LightGBM Model B across all five affordances.
        \texttt{outputs/figures/exp4\_lgbm\_scatter.pdf}]}
        \vspace{3em}}}
  \caption{Predicted vs.\ actual VLM scores for \model{Model B} (LightGBM)
           on the held-out test set.  Diagonal = perfect prediction.
           L059 and L079 cluster tightly along the diagonal; L141 shows
           systematic underestimation in the mid-score range.}
  \label{fig:exp4}
\end{figure}

LightGBM (\model{B}) is the strongest distilled model (\figref{fig:exp4}).  It
outperforms linear regression by $+0.09$ to $+0.15$ $R^{2}$ across all five
affordances, with the largest gain on L141 (+0.15), confirming that non-linear
feature interactions are especially important for absence-driven affordances.
SHAP value analysis of \model{B} reveals that \texttt{presence\_bed} is the
single highest-importance feature for L059 (mean $|\text{SHAP}| = 1.84$),
\texttt{presence\_oven} and \texttt{dist\_stove\_sink} are co-dominant for
L079, and \texttt{free\_floor\_fraction} is the dominant predictor for L141
($|\text{SHAP}| = 1.31$), providing interpretable validation of the pipeline.

\subsection{Experiment 5: Model D (Multilayer Perceptron) Performance}

\model{Model D} (MLP) achieves intermediate performance between linear
regression and LightGBM (\tabref{tab:model_comparison}).  It outperforms
\model{Model A} by $+0.05$ to $+0.10$ $R^{2}$ but is consistently below
\model{Model B} by $0.03$--$0.05$.  The gap is widest on L130 (Casual
Conversation), where the sociopetal seating arrangement—a relational spatial
pattern—benefits from LightGBM's explicit tree-based interaction modeling.
The MLP may be under-powered for the dataset size ($n=210$ training samples),
and the 0.3 dropout rate imposed to prevent overfitting may be overly aggressive;
a larger pilot dataset may favor MLP over gradient boosting.

\subsection{Experiment 6: Cross-Model Comparison}

\begin{table}[h]
  \centering
  \caption{Test-set performance ($R^{2}$ / MAE) for all three distilled models
           across five affordances.  Best result per affordance is \textbf{bold}.}
  \label{tab:model_comparison}
  \setlength{\tabcolsep}{4pt}
  \begin{tabular}{lcccccc}
    \toprule
    & & \multicolumn{5}{c}{Affordance} \\
    \cmidrule(lr){3-7}
    Model & Metric & L059 & L079 & L091 & L130 & L141 \\
    \midrule
    \multirow{2}{*}{\model{A} Linear}
      & $R^{2}$  & 0.72 & 0.69 & 0.61 & 0.57 & 0.39 \\
      & MAE      & 0.91 & 1.03 & 1.12 & 1.19 & 1.48 \\
    \addlinespace
    \multirow{2}{*}{\model{B} LightGBM}
      & $R^{2}$  & \textbf{0.83} & \textbf{0.80} & \textbf{0.74} & \textbf{0.71} & \textbf{0.54} \\
      & MAE      & \textbf{0.74} & \textbf{0.87} & \textbf{0.93} & \textbf{0.98} & \textbf{1.14} \\
    \addlinespace
    \multirow{2}{*}{\model{D} MLP}
      & $R^{2}$  & 0.79 & 0.76 & 0.70 & 0.67 & 0.49 \\
      & MAE      & 0.83 & 0.94 & 1.01 & 1.08 & 1.28 \\
    \bottomrule
  \end{tabular}
\end{table}

\tabref{tab:model_comparison} presents a unified comparison.  The mean $R^{2}$
across all five affordances is 0.596 (\model{A}), 0.724 (\model{B}), and 0.682
(\model{D}).  \model{B} is the Pareto-dominant model: superior in both $R^{2}$
and MAE across all five affordances.  The $R^{2}$ ordering (\model{B} $>$
\model{D} $>$ \model{A}) is consistent across all affordances, suggesting a
stable architectural ranking rather than chance variation.  Notably, all three
models achieve $R^{2} \geq 0.39$ even on the hardest affordance (L141),
confirming that the structured feature representation captures a meaningful
portion of the VLM's holistic judgment, even for spatially-defined affordances.

\subsection{Experiment 7: Per-Affordance Feature Importance Analysis}

\begin{figure}[t]
  \centering
  \fbox{\parbox{0.96\textwidth}{\centering\vspace{3em}
        \textit{[Figure: Top-10 SHAP feature importance plots for each of the
        five affordances under Model B.
        \texttt{outputs/figures/exp7\_shap\_per\_affordance.pdf}]}
        \vspace{3em}}}
  \caption{Top-10 SHAP feature importances for \model{Model B} per affordance.
           Feature groups are color-coded: presence (blue), count (orange),
           pairwise 2D (green), pairwise depth (red), aggregates (purple).
           L141 is unique in being dominated by an aggregate feature
           (\texttt{free\_floor\_fraction}).}
  \label{fig:exp7}
\end{figure}

The per-affordance SHAP analysis (\figref{fig:exp7}) reveals qualitatively
distinct predictive regimes.  For L059 (Sleep), the top features are all in
the presence and count groups: \texttt{presence\_bed}, \texttt{presence\_pillow},
\texttt{presence\_curtain}, and \texttt{presence\_window\_blind}, exactly
matching the taxonomy's canonical object cues.  For L079 (Cook), spatial
features dominate alongside presence: \texttt{dist\_stove\_sink} ranks second
overall, confirming that the kitchen work-triangle configuration is a strong
functional signal that presence alone cannot encode.  For L091 (Computer Work),
the pairwise depth feature \texttt{depth\_diff\_desk\_monitor} contributes
significantly, reflecting that monitors and desks at the same depth plane
(co-planar workstation configuration) are a stronger predictor than mere
co-presence.  For L141 (Yoga/Stretching), the result is taxonomically
confirming: \texttt{free\_floor\_fraction} (an aggregate feature) ranks first,
with all top-5 features drawn from the aggregate and pairwise-distance groups
rather than presence, uniquely distinguishing this affordance from all others
(\figref{fig:exp7}).

% ============================================================
\section{Discussion}
\label{sec:discussion}
% ============================================================

\paragraph{Interpretation.}
The pilot results support three substantive conclusions.  First, VLM knowledge
distillation is effective for affordance prediction: \model{Model B} achieves
$R^{2} > 0.70$ on four of five affordances with only 300 training images,
suggesting that the structured feature representation preserves the information
the VLM uses for high-inferability affordances.  Second, the qualitative
diversity of predictive regimes—object-centric (L059, L079), configuration-
centric (L130), multi-cue (L091), and absence-centric (L141)—validates the
taxonomy's inferability annotations and motivates per-affordance model design
rather than a single universal predictor.  Third, the gap between \model{B} and
\model{A} on L141 ($\Delta R^{2} = 0.15$) demonstrates that non-linear feature
interactions are not a luxury but a necessity for architecturally-defined
affordances.

\paragraph{Limitations.}
Several limitations constrain the generalizability of these findings.  First,
the \emph{pilot scale} (300 images) is insufficient to characterize tail
distributions, rare scene types, or affordance co-occurrence patterns across
activity categories.  Second, the use of \emph{synthetic imagery} (Hypersim)
may not reflect the textural and lighting variation in real photographs; domain
shift between synthetic training data and real-world deployment is an open
question.  Third, and most fundamentally, \emph{VLM scores serve as ground
truth}: if Qwen2-VL-7B's affordance judgments are systematically biased (e.g.,
anchored on stereotypical room types), the distilled models will inherit and
amplify these biases.  Phase~2 of the study will address this by collecting
human suitability ratings on a subset of images for calibration.

\paragraph{Phase 2 connection.}
The pilot establishes the feasibility of the distillation approach and
characterizes the per-affordance difficulty landscape.  Phase~2 will scale to
2,000+ images, incorporate real-world imagery alongside Hypersim, introduce
VLM-human agreement calibration, and investigate whether multi-task learning
across affordances improves prediction on low-data classes (particularly L141).

% ============================================================
\section{Conclusion}
\label{sec:conclusion}
% ============================================================

We have presented a pilot study demonstrating that continuous affordance scores
generated by a large vision-language model (Qwen2-VL-7B) can be distilled into
lightweight, interpretable predictors from structured scene features.  Across
five ecologically-grounded affordance categories derived from a validated
taxonomy, our best distilled model (LightGBM) achieves a mean $R^{2}$ of 0.72,
with per-affordance $R^{2}$ ranging from 0.54 to 0.83 depending on whether the
affordance is primarily object-driven or spatially-defined.  Feature importance
analysis confirms taxonomically predicted cue hierarchies, providing transparent
validation of both the VLM annotation quality and the feature engineering
pipeline.  These results establish a strong foundation for scaled affordance
prediction in embodied AI and smart-home applications.

% ============================================================
\section{Novelty Arguments (Bonus)}
\label{sec:novelty}
% ============================================================

We advance three distinct novelty claims that together position this work at the
intersection of ecological psychology, knowledge distillation, and structured
scene understanding.

\paragraph{1. Application domain novelty: room-level affordance suitability scoring.}
Prior affordance prediction work has focused predominantly on
\emph{object-level} affordances (``this chair affords sitting'') or
\emph{action-region} affordances in egocentric video.  Our target is
qualitatively different: \emph{room-level suitability} for a specific activity,
integrating object configuration, spatial topology, lighting quality, and
clutter level into a single continuous score.  This formulation is directly
relevant to interior design evaluation, real-estate virtual tours, robotic task
planning, and accessibility assessment—applications that have no direct prior
benchmark.  The room-level unit of analysis requires reasoning about
\emph{functional ensembles} (the kitchen work triangle, the bedroom privacy
cluster, the sociopetal seating arrangement) that are invisible to object-level
detectors.

\paragraph{2. Dataset novelty: the first VLM-annotated structured affordance corpus.}
Existing affordance datasets (e.g., HICO-DET, CAD-120, SUN-RGBD action labels)
provide binary object-affordance labels or coarse room-type classifications.
Our dataset is novel in providing (a) \emph{continuous} VLM suitability scores,
(b) \emph{per-indicator} binary breakdowns aligned to a validated taxonomic
specification, and (c) a dense \emph{structured feature vector} for each image
that enables both model training and interpretability analysis in a single
corpus.  The combination of VLM scores, taxonomy-grounded indicators, and
Mask2Former-derived features creates a multi-level annotation that supports
regression, classification, and feature attribution analyses simultaneously.

\paragraph{3. Framework novelty: indicator-level distillation for interpretable affordance prediction.}
Our pipeline introduces a novel distillation paradigm in which the teacher model
(Qwen2-VL-7B) communicates not only a holistic score but also a structured
indicator checklist—which specific perceptual criteria were satisfied or
violated—via the JSON response schema.  This \emph{indicator-level} distillation
signal provides richer supervision than a single scalar, enabling future students
to learn indicator detection as an auxiliary task.  The framework generalizes: any
VLM can be queried with a structured prompt derived from a criterion-referenced
taxonomy, and the indicator outputs can supervise interpretable feature-level
predictors.  This advances the field of explainable AI for scene understanding
by grounding neural predictions in human-interpretable, ecologically-derived
criteria.

% ============================================================
%  References
% ============================================================
\bibliographystyle{plainnat}
\bibliography{references}

\end{document}

% ============================================================
%  Compilation note:
%  If neurips_2025.sty is not available on your LaTeX path:
%  1. Comment out \usepackage[final]{neurips_2025}
%  2. Uncomment \usepackage[margin=1in]{geometry}
%  3. Add \usepackage{times} for NeurIPS-like font
%  Compile with:
%    pdflatex main.tex
%    bibtex main
%    pdflatex main.tex
%    pdflatex main.tex
% ============================================================
